../cictikz/src/cictikz/data/tex/cictikz_preamble.tex

% The Active-RC integrator the lecture uses to talk about what a real OTA
% does to it: R from V_I to the virtual ground, C in feedback, + grounded.
% R and C carry no labels in the original -- the slide is about A_0 and
% omega_ta, not about component values.

\begin{circuitikz}[american, thick, transform shape]

  ../cictikz/src/cictikz/data/tex/cictikz_lib.tex

  \def\yin{1.8}   % inverting input, and the R branch
  \def\yfb{3.2}   % the feedback capacitor

  % R from the input to the virtual ground
  \node[anchor=east] at (0.3,\yin) {$V_{I}$};
  \draw (0.4,\yin) -- (1.0,\yin);
  \draw (1.0,\yin) \hresistor{};
  \draw (resEnd) -- (3.2,\yin);

  % OTA, inverting input on top
  \draw (3.2,\yin) \cicOtaSSWP{$-$}{$+$};
  \draw (cicOtaS_inn) -- (2.35,{\yin-0.8}) -- (2.35,0.6);
  \draw (2.35,0.6) \vground;

  % C in feedback
  \draw (2.9,\yin) -- (2.9,\yfb);
  \fill (2.9,\yin) circle (0.07);
  \draw (2.9,\yfb) -- ++(0.85,0) \hcapacitor{};
  \draw (capEnd) -- (6.2,\yfb) -- (6.2,{\yin-0.4});

  % output
  \draw (cicOtaS_out) -- (6.9,{\yin-0.4});
  \fill (6.2,{\yin-0.4}) circle (0.07);
  \node[anchor=west] at (6.95,{\yin-0.4}) {$V_{O}$};

\end{circuitikz}
\end{document}
